% ===== PRÉAMBULE CANONIQUE — à coller VERBATIM en tête de CHAQUE .tex =====
\documentclass[11pt,a4paper]{article}
\usepackage[utf8]{inputenc}\usepackage[T1]{fontenc}\usepackage{lmodern}
\usepackage[french]{babel}
\usepackage{amsmath,amssymb,mathtools}\usepackage{icomma}
\usepackage[version=4]{mhchem}          % \ce{...} : équations & formules
\usepackage{siunitx}\sisetup{locale=FR} % \qty{}{}, \si{}, \ang{}
\usepackage{chemfig}                    % \chemfig{...} : structures développées
\usepackage{xcolor}\usepackage{colortbl}
\usepackage{tikz}\usepackage{pgfplots}\pgfplotsset{compat=1.18}
\usepackage[most]{tcolorbox}\usepackage{enumitem}\usepackage{array,booktabs}
\usepackage[a4paper,margin=2.2cm]{geometry}
\usetikzlibrary{arrows.meta,calc,angles,quotes,positioning,patterns,backgrounds,shapes.geometric}
\definecolor{primary}{RGB}{222,49,99}\definecolor{primarydark}{RGB}{150,22,55}
\definecolor{defcol}{RGB}{0,114,178}\definecolor{methcol}{RGB}{52,92,148}
\definecolor{excol}{RGB}{0,135,90}\definecolor{attncol}{RGB}{200,40,55}
\tcbset{boxbase/.style={enhanced,breakable,boxrule=0.9pt,arc=2.5pt,
  left=3mm,right=3mm,top=2.3mm,bottom=2.3mm,fonttitle=\bfseries,coltitle=white}}
% --- boîtes TUTORIEL (noms EXACTS, ne pas renommer) ---
\newtcolorbox{cle}[1][]{boxbase,colback=defcol!6,colframe=defcol,
  title={\textsc{À retenir}\ifx\relax#1\relax\relax\else\textmd{ : \itshape #1}\fi}}
\newtcolorbox{methode}[1][]{boxbase,colback=methcol!7,colframe=methcol,sharp corners,
  title={\textsc{Méthode}\ifx\relax#1\relax\relax\else\textmd{ --- \itshape #1}\fi}}
\newtcolorbox{exemplebox}[1][]{boxbase,colback=excol!5,colframe=excol,
  title={\textsc{Exemple}\ifx\relax#1\relax\relax\else\textmd{ : \itshape #1}\fi}}
\newtcolorbox{piege}[1][]{boxbase,colback=attncol!6,colframe=attncol,
  title={\textsc{Piège}\ifx\relax#1\relax\relax\else\textmd{ : \itshape #1}\fi}}
\newtcolorbox{remarque}[1][]{boxbase,colback=black!4,colframe=black!45,
  title={\textsc{Remarque}\ifx\relax#1\relax\relax\else\textmd{ : \itshape #1}\fi}}
% --- boîtes EXERCICE / CORRIGÉ (noms EXACTS, auto-comptées) ---
\newtcolorbox[auto counter]{exo}[1][]{boxbase,colback=primary!4,colframe=primary,
  title={\textsc{Exercice}~\thetcbcounter\ifx\relax#1\relax\relax\else\textmd{ --- \itshape #1}\fi}}
\newtcolorbox[auto counter]{corrige}[1][]{boxbase,colback=excol!4,colframe=excol,
  title={\textsc{Corrigé}~\thetcbcounter\ifx\relax#1\relax\relax\else\textmd{ --- \itshape #1}\fi}}
\newcommand{\R}{\mathbb{R}}\newcommand{\N}{\mathbb{N}}\newcommand{\Z}{\mathbb{Z}}\newcommand{\ds}{\displaystyle}
% (un \usepackage{fancyhdr}+en-tête est possible ; il est ignoré côté web)
% =========================================================================
\begin{document}
\sisetup{per-mode=symbol}
\begin{center}\small\itshape On rappelle : $\log 2 \approx 0{,}30$,\quad $\log 3 \approx 0{,}48$,\quad $\log 5 \approx 0{,}70$.\end{center}

\begin{exo}[pH non entier, dans les deux sens]
\begin{enumerate}[label=\alph*)]
  \item Une solution a $[\ce{H3O+}] = \qty{2,0e-5}{mol\per L}$. Calcule son pH.
  \item Une solution a $\mathrm{pH} = 3{,}0 + \log 2$. Exprime $[\ce{H3O+}]$ sous la forme $a\times10^{n}$.
  \item Pour la solution du a), calcule aussi $[\ce{OH-}]$ (à \qty{25}{\celsius}).
\end{enumerate}
\end{exo}

\begin{exo}[passer de $\mathrm{p}K_a$ à $\mathrm{p}K_b$]
Le couple \ce{NH4+/NH3} a $\mathrm{p}K_a = 9{,}2$ (à \qty{25}{\celsius}).
\begin{enumerate}[label=\alph*)]
  \item Calcule le $\mathrm{p}K_b$ de la base \ce{NH3}.
  \item En déduis $K_b$ de \ce{NH3} sous la forme $a\times10^{n}$.
  \item L'ammoniac est-il une base forte ou faible ? Justifie.
\end{enumerate}
\end{exo}

\begin{exo}[force croisée d'un couple]
On donne trois acides et leur $\mathrm{p}K_a$ : \ce{HNO2} ($3{,}3$), \ce{CH3COOH} ($4{,}8$), \ce{HClO} ($7{,}5$).
\begin{enumerate}[label=\alph*)]
  \item Range les \textbf{bases conjuguées} \ce{NO2-}, \ce{CH3COO-}, \ce{ClO-} de la plus \textbf{forte} à la
        plus faible.
  \item Justifie la règle : « plus un acide est faible, plus sa base conjuguée est forte ».
\end{enumerate}
\end{exo}

\begin{exo}[par le pOH]
Une solution basique a $[\ce{OH-}] = \qty{5,0e-3}{mol\per L}$ (à \qty{25}{\celsius}).
\begin{enumerate}[label=\alph*)]
  \item Calcule le pOH.
  \item En déduis le pH.
  \item Vérifie le caractère (acide / neutre / basique) attendu.
\end{enumerate}
\end{exo}

\begin{exo}[quand la température change le pH neutre]
On chauffe de l'eau pure jusqu'à une température $T$ où $K_e = \qty{4,0e-14}{}$.
\begin{enumerate}[label=\alph*)]
  \item Dans l'eau pure, $[\ce{H3O+}] = [\ce{OH-}]$. Calcule $[\ce{H3O+}]$ à cette température.
  \item En déduis le pH neutre à $T$.
  \item Une solution de pH $= 7$ à cette température est-elle acide, neutre ou basique ?
\end{enumerate}
\end{exo}

\end{document}
